\section{Materials and Methods}
\label{sec:methods}

\subsection{Problem definition}

Let $y_{1:T}$ be a univariate series with seasonal period $m$. An \emph{instance} is a
triple $(y, t, h)$: forecast $h$ steps from origin $t$ using only $y_{1:t}$. A base
forecaster $f$ produces a point forecast $\hat{y}_{t+1:t+h}$ and a predictive
distribution. The agent may then apply a corrective action $a$ from a finite set
$\mathcal{A}$, obtaining a revised forecast.

The object of study is the \emph{correction policy} $\pi$, mapping evidence available at
time $t$ to an action in $\mathcal{A}$. Writing $L(a)$ for the realised loss under
action $a$, we evaluate policies by
\begin{equation}
  \mathrm{Regret}(\pi) = \mathbb{E}\big[L(\pi(\cdot)) - \min_{a \in \mathcal{A}} L(a)\big],
\end{equation}
which scores the \emph{choice} instead of the forecast, and is computable only because
the benchmark stores $L(a)$ for every $a$.

\subsection{TS-AgentBench}

\subsubsection{Corrective action space}
\label{sec:actionspace}

Ten typed actions, each paired with the failure modes it is designed to repair:
accept (the null action); switch model class; refit on the segment after the most recent
structural break; robustify by replacing outliers with local medians; damp the
extrapolated trend; re-specify the seasonal period; apply a variance-stabilising
transform; ensemble the best-backtesting models; recalibrate prediction intervals; and
abstain by falling back to a seasonal-naive forecast with a low-confidence flag.

The space is deliberately small, closed and typed. Typing is what makes diagnosis
consequential: the agent does not merely observe that something is wrong, it commits to
a failure mode and applies the matched repair, and the benchmark records whether that
was correct. Interval recalibration is retained even though it cannot change point
accuracy, because it separates ``the forecast was wrong'' from ``the forecast was
acceptable but over-confident''---a distinction invisible to point metrics.

\subsubsection{Counterfactual outcome tensor}

For every instance we execute all ten actions and score each against the realised
target, storing point metrics (MAE, RMSE, MASE \citep{hyndman2006mase}, sMAPE, WAPE) and
probabilistic metrics (CRPS, interval coverage, width, Winkler score). MASE is the
primary point metric because it is scale-free and defined for series containing zeros,
and its denominator is computed on the training history only; scaling by the full
series would leak test information into every reported number.

The primary probabilistic metric is CRPS, computed from sample paths in its energy form
so that it is available for models with no closed-form predictive distribution. The
choice matters for the interpretation of our results rather than only for convenience:
CRPS is a \emph{strictly proper} scoring rule \citep{gneiting2007crps}, minimised
uniquely by the true predictive distribution. A policy therefore cannot improve CRPS by
inflating interval widths until coverage looks acceptable, which is the obvious
alternative explanation for any claim that correction improves probabilistic accuracy.
We report coverage and Winkler scores alongside it as mechanism, never as the outcome,
and Section~\ref{sec:overconfidence} shows empirically what widening costs on this
benchmark. From these we derive per instance:
the oracle action and its loss; a \texttt{should\_correct} label, true when some action
beats accepting by more than a relative margin $\delta$; the number of actions that
would have caused harm; and the regret of any candidate policy.

Two consequences follow. Evaluating a new correction policy requires no model fitting,
so the full ablation and robustness grid is affordable on a CPU. And because the
stochastic work happens once, replays are exactly reproducible.

\paragraph{Selection bias, and why we report two oracles.}
The unrestricted oracle takes a minimum over ten noisy outcomes on a single
realisation. The minimum of ten draws is biased downward, so part of the apparent
headroom is the winner's curse instead of a repair anyone could have chosen. We
therefore also report a \emph{realizable oracle}: the best action selected on the
\emph{validation} origins of the same series and then applied to test. It involves no
selection on the evaluation outcome and is attainable in deployment. The difference
between the two is the selection bias, and we report it explicitly
(Section~\ref{sec:results}); an evaluation quoting only the unrestricted oracle
overstates what any policy could achieve.

\subsubsection{Data}

\emph{Synthetic.} Fourteen generator families supply ground truth that real data cannot:
linear and damped trends, single and multiple seasonality, stationary AR and
near-unit-root dynamics, level shifts, gradual drift, additive and innovational
outliers, variance shifts, Markov regime switching, short histories, missing blocks,
and multivariate VAR with cross-series dependence. Each series records its true change
points, outlier indices, regime labels and expected failure mode. Structural breaks are
positioned in a late window, straddling the validation/test boundary, because whether a
break precedes or follows the forecast origin determines whether \emph{any} correction
could help---and conditioning on that is what H2 requires.

\emph{Real.} Public datasets with documented licences and recorded checksums, drawn from
the Monash archive (CC~BY~4.0) \citep{monash2021} and the ETT collection
\citep{informer2021}. Provenance for every file---source URL, licence, SHA-256, access
date---is generated by the preparation script into \texttt{datasets/LICENSES.md}. Any
dataset whose licence we could not verify was dropped rather than assumed. Real series
carry no ground-truth structural labels, and the analysis conditions on that flag rather
than treating detector output as truth.

\subsubsection{Splits and leakage control}

Splits are strictly temporal ($\text{train} < \text{val} < \text{test}$) with
rolling-origin evaluation. Scalers, feature statistics and MASE denominators are
computed on training data only. The reliability estimator and every threshold are fitted
on validation only. Oracle labels derive from test outcomes and are therefore
evaluation-only artifacts: they live in a module that agent code cannot import, a
constraint enforced by a test that walks the import graph and fails the build on
violation. Further tests corrupt all post-origin observations and assert that no
diagnostic or corrective action changes its output, which catches leakage that a
signature inspection would miss.

\subsection{Forecasting models}

The pool spans naive, seasonal naive, drift, historical mean, Theta, ETS, ARIMA with a
compact AIC-scored order search, ridge autoregression, random forest, gradient boosting,
an LSTM, a temporal convolutional network, and a patch Transformer in the style of
\citet{patchtst2023}. Tree models are
fitted on first differences and integrated back, which is standard practice and avoids a
straw-man baseline: undifferenced tree ensembles cannot extrapolate and would flatline
on any trending series. Neural models emit all $h$ steps at once with instance
normalisation, and share an identical training loop so that differences are attributable
to architecture rather than to tuning.

Pretrained time-series foundation models \citep{chronos2024} are supported by the
released code but are not part of the pool reported here: running one across the full
counterfactual tensor exceeds our CPU budget, and including it on a subset would make
the action space differ between instances, which the decomposition analysis relies on it
not doing. This is a stated limitation instead of a claim that the finding is
architecture-independent.

All models expose the same probabilistic interface, producing sample paths from which
intervals and CRPS are derived. Uniform probabilistic output is what allows a naive
baseline and a neural network to be compared on CRPS at all; without it, probabilistic
comparison silently becomes a comparison of which models had an analytic interval
implemented.

\paragraph{Controlled base-model selection.}
Every arm shares one base-model selection procedure: the model with the lowest
rolling-origin backtest error on history. Arms differ \emph{only} in their correction
policy. Allowing each arm to also choose its own model would confound the variable under
study, since a difference in final accuracy could then arise from better model selection
rather than better correction decisions.

\subsection{Diagnostics}

Eight history-only tools supply the evidence: summary features, STL decomposition with
trend and seasonal strength, ACF/PACF, seasonal-period detection, ADF and KPSS
stationarity tests, PELT change-point detection, robust anomaly detection on
decomposition remainders, and distribution-shift statistics comparing a recent window
against earlier history. Each declares a computational cost, which is what makes H4's
cost claim measurable rather than asserted.

Two implementation details were determined empirically against families whose structure
is known by construction. Change-point detection operates on a standardised, linearly
detrended series with an $\ell_2$ cost: the default kernel cost is unusable on
unstandardised data, returning nothing at all on a textbook level shift when the series
level is large, and an untreated linear trend is otherwise decomposed into a staircase
of spurious breaks. Anomaly detection uses median and MAD rather than mean and standard
deviation, because the latter are themselves corrupted by the outliers being sought.

\subsection{Reliability-gated correction}

RGC proceeds in four stages.

\paragraph{1. Gate.}
A gradient-boosted classifier, fitted on validation instances and isotonically
calibrated, predicts $P(\text{some repair beats accepting})$. Correction is triggered
when this exceeds a threshold chosen on validation to meet a correction budget. We gate
on predicted \emph{gain} rather than predicted \emph{error} deliberately: a forecast can
be doomed and unrepairable---for instance when a structural break falls inside the
forecast horizon---and spending computation there is waste. A second head predicts
$P(\text{large error})$ and drives abstention and the calibration analysis.

Budget-first thresholding also makes the compute-matched comparison exact: the
rate-matched random arm corrects the same number of instances, so the two arms differ
only in \emph{which} instances they choose.

\paragraph{2. Diagnose.}
Signals are mapped to a failure mode by transparent rules, restricting candidates to
matched repairs. The rules are deliberately not learned: a learned diagnoser would
likely score better but would make the ablations uninterpretable, since removing a
signal family could be silently compensated by the model rebalancing. We evaluate an LLM
router as an alternative to these rules, treating it as a comparison instead of an
assumption.

\paragraph{3. Validate.}
Each candidate is scored on the validation origins of the same series, and a correction
is adopted only if it improves on accepting by more than $\delta$. This is the crux of
the method. Those origins lie strictly in the past relative to the test origin, so a
deployed agent could compute exactly this by backtesting candidate repairs on recent
history. The agent never judges its own output; it measures. This is the concrete answer
to the finding that self-correction fails when a model is left to assess itself
\citep{huang2024selfcorrect,kamoi2024selfcorrect}.

\paragraph{4. Abstain.}
If predicted risk remains high after correction, the agent falls back to a seasonal-naive
forecast and flags low confidence. Abstention in forecasting cannot mean returning
nothing---downstream systems still require a number---so it takes the form of retreating
to the most defensible simple forecast while signalling that the estimate should not be
trusted.

\subsection{Experimental arms}
\label{sec:arms}

Every policy reported anywhere in the paper is listed here, marked as pre-registered or
added later. Later additions are recorded as deviations in the analysis plan.

\paragraph{Pre-registered.}
\begin{description}
  \item[\texttt{never\_correct}] Accept the base forecast. The reference every other
  arm is compared against.
  \item[\texttt{always\_correct}] Correct every instance with the best-validated repair.
  \item[\texttt{random\_correct}] Correct a random subset at the rate the gate fires.
  The compute control for H6.
  \item[\texttt{threshold\_single}] Correct when one signal, backtest error, exceeds a
  validation-fitted threshold.
  \item[\texttt{best\_fixed}] One repair per series, chosen on validation: the
  realizable oracle. Its result bears directly on Section~\ref{sec:decomposition}.
  \item[\texttt{rgc}] The reliability-gated policy as pre-registered, with a fixed 2\%
  acceptance margin and a target correction rate of 20\%.
  \item[\texttt{llm\_single}, \texttt{llm\_critique\_gate}] A language model that
  chooses a repair in one pass, or first decides by self-critique whether to repair.
\end{description}

\paragraph{Added later.}
\begin{description}
  \item[\texttt{rgc\_calibrated}] A revised gate with a per-action gain calibrator
  (deviation~D4). Its threshold is fitted separately for each metric, so it acts on very
  few instances under MASE and on nearly all under CRPS, where it mostly recalibrates
  intervals.
  \item[\texttt{always\_ensemble}] Apply the ensemble action everywhere. Chosen
  after the main results; Section~\ref{sec:staticcheck} shows it is the action
  validation selects on CRPS, but not on MASE.
  \item[\texttt{rgc\_llm\_router}] \texttt{rgc}'s gate decides whether to correct; a
  language model chooses which repair (deviation~D8).
\end{description}

\paragraph{Reading the tables.}
\texttt{rgc} and \texttt{rgc\_calibrated} are different policies and should not be
compared across tables as if they were one. Their numbers also differ by aggregation:
Table~\ref{tab:headline} reports series-level means, and the decision-quality and cost
tables report instance-level means. The LLM tables use a smaller subsample, one
forecast origin per series, so none of their levels match the headline table's.

\subsection{LLM configuration and contamination control}

LLM arms run at temperature 0 on three model families: \texttt{gemini-3.5-flash-lite},
a small hosted proprietary model; \texttt{qwen3.8-27b}, an open-weight model of
different lineage served by a different provider; and \texttt{gpt-oss-20b}, an
open-weight reasoning model run at \texttt{reasoning\_effort} \texttt{= medium}
with a completion budget sized for the reasoning trace as well as the answer. Free-tier
quotas, which meter requests per day for one provider and tokens per day for another,
set how much of the benchmark each family covers; coverage is reported per row rather
than averaged over, and every non-LLM arm is re-scored inside each family's own
subsample so that rows are comparable within a family and not across families.

A fourth condition tests whether the zero-shot setting is doing the work. The LLM arms
are zero-shot while the arithmetic estimator they are compared against is fitted on the
validation split, so a null result would otherwise be ambiguous between a model that
cannot read the evidence and one that was never shown what a good decision looks like.
In that condition eight worked examples, drawn from validation only and labelled by the
counterfactual tensor, are prepended to every prompt. Their class mix follows the
validation base rate instead of being balanced --- a balanced prefix would itself hint
at how often to intervene, which is the quantity under study --- and the positive
examples are spread across distinct oracle actions so the prefix does not teach a single
repair. The example set is fixed across instances and across metrics, so the only thing
differing between two prompts is the instance being asked about, and the arm makes one
decision per instance that is then scored under both metrics, exactly as the zero-shot
arms do.

Every response is cached under a hash of the provider, model, temperature, maximum
output tokens, system prompt, user prompt and, where set, the reasoning effort. Reruns
therefore replay identical responses and make no network calls: every LLM result in this
paper reproduces without an API key or a provider account. Calls, cache hits, tokens and
latency are counted, because a cost claim that is not measured is not evidence.

One failure mode needed an explicit guard. A malformed or truncated reply falls back to
accepting the forecast, and on a reasoning model truncation usually means the completion
budget was spent on the chain of thought instead of the answer. Left unchecked that
would bias the intervention rate downward, which is the single quantity these findings
rest on. Provider failures and unparsed replies are therefore counted together against a
2\% ceiling and reported separately; a run that breaches it is aborted rather than
reported. Every arm in this paper came in at 0\% on both counts.

Several real datasets are long-standing and public, so memorisation is a live risk.
Prompts therefore omit the dataset name, series identifier, all dates and all raw
values, presenting only scale-free diagnostic statistics: the model is asked to reason
about a diagnosis, never to recall a series. Causal claims (H2) rest on freshly seeded
synthetic series that cannot have been memorised.

\subsection{Statistical analysis}

All confirmatory tests were fixed in advance. Paired comparisons across series use the
Wilcoxon signed-rank test with Cliff's delta as effect size; forecast-accuracy pairs use
Diebold--Mariano \citep{diebold1995} with the Harvey--Leybourne--Newbold small-sample
correction \citep{harvey1997}; confidence intervals are non-parametric bootstrap with
\textbf{resampling at the series level}.
That last choice is consequential rather than cosmetic: instances from one series share
a model fit, a history and a regime, so resampling instances would shrink every interval
and manufacture significance. It is also the choice we initially failed to apply; deviation~D7 in the analysis plan records that the first draft reported instance-level
$p$-values, and that correcting this removed two claims. Equivalence claims (H4) use two
one-sided tests, since an absent difference is not evidence of equivalence. The seven
confirmatory hypotheses form one family under Holm--Bonferroni correction
\citep{holm1979}; per-stratum and per-horizon breakdowns are exploratory and labelled as
such.

Where a rank test and a bootstrap interval disagree --- which happens here, because the
loss is heavy-tailed --- we report both instead of the more favourable one. They answer
different questions: the rank test asks whether a policy wins more often than it loses,
the interval asks whether the sizes of those wins establish a mean improvement. A policy
can do the first and not the second, and describing such a policy as simply ``better''
would overstate what the data shows.
